\section{Filtered local systems on surfaces}
\label{sec:diagrams}

The proof uses auxiliary surfaces whose boundary local systems are related
by filtrations.  Cutting these surfaces and refining their filtrations
will reduce the rank.  We first define the auxiliary objects and the
linear algebra needed to reconstruct them after a cut.

Use the field $F$ and its full ring of integers $R$ from
Section~\ref{sec:arithmetic}.  All prescribed
scalars below are roots of unity in $F$.  A local system over $R$ is
required to have finite free fibers.  Its parallel transports are thus
isomorphisms over the full ring $R$.

\subsection{Surfaces, seams, and their solutions}

A \emph{surface diagram} has finitely many compact oriented surfaces,
called its facets.  Each facet has an assigned rank between $0$ and $r$
and finitely many interior punctures.  A scalar monodromy is prescribed
at each puncture, using the positive orientation around that puncture.

Every boundary circle is either left unmarked or divided into intervals
by finitely many marked points.  These boundary pieces are grouped into
\emph{seams}.  A seam has a distinguished parent side and an ordered
list of child sides, written
\[
 E:(W_1,\ldots,W_m),\qquad
 \rk E=\sum_{j=1}^m\rk W_j.
\]
All its sides are parametrized by the same interval or circle.  The
boundary orientation of each child is opposite to that of the parent.
Different side occurrences may belong to the same facet.  An unmarked
circle seam uses entire circles and has no permutation of its sides.

At an interval endpoint, the boundary of each facet pairs its two
incident side germs.  We permit precisely the following vertices.
\begin{enumerate}
\item A parent list continues with the same ordered children, or with
two adjacent children interchanged.  Each child continues to itself.
\item A child $W$ of one parent list is replaced by an ordered list
$W'_1,\ldots,W'_s$, and there is a third seam
$W:(W'_1,\ldots,W'_s)$.  The parent germs continue to each other;
the two occurrences of $W$ are paired; the two occurrences of each
$W'_j$ are paired; and unchanged children continue to themselves.
The reverse operation is also permitted.
\end{enumerate}
The specification concerns distinct occurrences of germs, even when
some belong to one facet.  No embedding of the diagram in a surrounding
space is required.

A \emph{solution} assigns to each punctured facet a local system of its
rank, with the prescribed scalar monodromies.  At a seam it assigns a
parallel filtration of the parent and parallel identified quotients
\[
 0=E_0\subset E_1\subset\cdots\subset E_m=E,
 \qquad E_j/E_{j-1}\simeq W_j.
\]
The conditions at vertices are as follows.  At a continuation, all data
continue.  At a split or merge, the finer parent filtration is obtained
by refining the $W$-quotient with its given filtration, and the
identified quotients agree.  At an interchange of $A,B$, the coarsened
filtration with quotient $A\oplus B$ is unchanged.  Within that quotient
the two intermediate subspaces are complementary, and the maps to
the identified quotients are the natural projections.  Thus the
interchange includes a specified direct-sum decomposition, not just a
pair of transverse flags.

Choose one frame on each connected facet.  Holonomies along a finite
set of paths, the seam flags, and their identified quotients describe
the solutions by finitely many algebraic equations and open conditions.
They form a complex variety, denoted $\mathcal S(D)$ for a diagram $D$.
Changing the auxiliary paths gives the same data through composition
and inversion of transport matrices.

An \emph{integral solution} uses free $R$-local systems and filtrations
whose identified graded modules are the specified free child fibers.
Complementarity at an interchange means a direct sum over $R$.
In particular, each filtration splits on a fiber: successive free
quotients are projective.  This definition does not assert that every
submodule of a free $R$-module is free.

\begin{theorem}[Diagram density]\label{thm:diagram-density}
For every surface diagram $D$ with ranks at most $r$ and prescribed
scalar puncture monodromies in $F$, its integral solutions over $R$
are Zariski dense in $\mathcal S(D)$.
\end{theorem}

We prove Theorem~\ref{thm:diagram-density} by induction on the largest
rank, normalizing identity seams before separating closed facets.
The nonclosed step is proved in Section~\ref{sec:reduction}.
The rank-one and genus-zero closed cases are elementary, and the
genus-one case is Lemma~\ref{lem:torus-density}.  At each remaining
rank, Lemma~\ref{lem:closed-density} gives the closed-surface step;
it uses the normalized nonclosed step at that same rank and the
completed induction at smaller ranks.

\subsection{Transferring density through choices}

\begin{lemma}\label{lem:density-transfer}
Let $Y$ be a complex variety.  For each member $\tau$ of a finite index
set, suppose there are morphisms
\[
 B_\tau\xleftarrow{\ p_\tau\ }T_\tau
 \xrightarrow{\ f_\tau\ }Y,
\]
where $p_\tau$ is open and surjective and the images of the $f_\tau$
cover $Y$.  Suppose that a set $I_\tau\subset B_\tau$ is dense and that,
over every point of $I_\tau$, a dense set of points in the fiber of
$p_\tau$ maps into a subset $J\subset Y$.  Then $J$ is dense in $Y$.
\end{lemma}

\begin{proof}
A nonempty open subset $O\subset Y$ has nonempty inverse image under
some $f_\tau$.  Its image under $p_\tau$ is nonempty and open, hence
contains a point of $I_\tau$.  The dense specified subset of that fiber
meets $f_\tau^{-1}(O)$.
\end{proof}

We apply this lemma with $B_\tau$ the solutions of a simpler diagram
and $T_\tau$ the choices for reconstruction.  The index $\tau$ records
ranks of finitely many intersections and projection images.  Its
finiteness follows from the fixed rank bound.  The forward reduction
need not be a morphism on the entire original solution variety; it is
enough that every original solution has a reconstruction in one of the
finitely many spaces $T_\tau$.

The choice spaces used below are locally affine bundles, or products
of these with general linear groups.  Their projections are therefore
open.  Over an integral solution, the stated free splittings identify
their parameters with matrices over $R$ and with invertible matrices
over $R$.  Both sets are Zariski dense in the corresponding complex
spaces by Lemma~\ref{lem:integral-parameters}.
We may include arbitrary frame changes among the choices;
coverage up to isomorphism then suffices for framed coverage.

\subsection{Comparing two flags}

The following lemma includes the compatibility of the identified
quotients.  That compatibility is needed when a cut surface is sewn
back together.
Flag decompositions and braid operations also underlie the geometric
stratifications of Mellit~\cite[Theorem~1.2]{Mellit} and
Su~\cite[Theorem~0.5]{Su}.
Here we formulate reconstruction over one ring of integers with
specified free graded modules, so that every change of splitting
has integral matrix parameters.

\begin{lemma}[Two-flag reconstruction]\label{lem:two-flags}
Let $F_\bullet,H_\bullet$ be two flags in a complex vector space $V$.
There are nonnegative integers $m_{ij}$ and a sequence of the allowed
split, interchange, and merge operations with these properties:
\begin{enumerate}
\item It refines $F_\bullet$ into blocks of ranks $m_{ij}$ in row-first
order, sorts them into column-first order by adjacent interchanges,
and merges them to $H_\bullet$.
\item Each inverted pair of blocks is interchanged once.  The flags
and all identified quotients in any such sequence admit a common
splitting into the labeled blocks.
\item Given two such pairs of flags, prescribed isomorphisms on their
$F$-graded modules and their $H$-graded modules extend to a full
isomorphism if they preserve the induced refinements and agree on
the double-graded modules.  The extensions form an affine space.
\end{enumerate}
The last two assertions hold over $R$ when the labeled graded modules
are free.  The affine spaces then have free integral parameters.
In families with fixed ranks, the spaces of extensions are locally
affine bundles.
\end{lemma}

\begin{proof}
Set
\begin{align*}
m_{ij}={}&\dim(F_i\cap H_j)-\dim(F_{i-1}\cap H_j)\\
 &-\dim(F_i\cap H_{j-1})+\dim(F_{i-1}\cap H_{j-1}).
\end{align*}
Split each intersection modulo the sum of the preceding row and column
intersections.  This gives
\[
 V=\bigoplus_{i,j}V_{ij},\qquad
 F_i=\bigoplus_{k\leq i,j}V_{kj},\qquad
 H_j=\bigoplus_{i,\ell\leq j}V_{i\ell}.
\]
Use these summands for the refinements.  Sort row-first into
column-first order, interchanging each inverted pair exactly once.
Every interchange is a direct-sum interchange.

Conversely, start by splitting the first fully refined flag of a
sequence as in the statement.  At an interchange of adjacent labels
$a,b$, the new $b$-subspace in the two-block quotient is the graph of
a map from the former $b$-summand to the $a$-summand.  Modify the
$b$-summand by that graph.  The label $a$ preceded $b$ in every earlier
order, since this pair has not previously crossed.  Adding an
$a$-component to the $b$-summand therefore preserves every earlier
prefix and its identified quotients.  The modified splitting realizes
the new flag and the natural quotient identifications at the
interchange.  Induction produces one splitting of all the flags.
The argument uses only splittings of free quotients and linear maps,
so also works over $R$.

For the extension assertion use common splittings on source and target.
A matrix preserving both flags has a block from $(i,j)$ to $(i',j')$
only if $i'\leq i$ and $j'\leq j$.  Its blocks in equal rows are
prescribed by the $F$-graded maps; those in equal columns are prescribed
by the $H$-graded maps.  The prescriptions coincide on their overlap
precisely when they agree on the double-graded modules.  All blocks
strictly lowering both indices are arbitrary.  The diagonal blocks
are isomorphisms, so the resulting matrix is invertible.  It also
respects every intermediate flag in the gallery.  Indeed each sorting
order extends the coordinatewise partial order on the labels $(i,j)$:
only incomparable labels are interchanged.  A block lowering the
indices therefore maps into an earlier prefix, and the diagonal blocks
give the prescribed maps on the labeled quotients.  This gives
the claimed affine space and its integral parameters.  For families,
split the constant-rank subbundles locally and apply the same block
description.
\end{proof}

\subsection{Lifting a filtered exact sequence}

Replacing a local system by a subspace $U$ and quotient $Q$ loses both
the position of a boundary flag in $U\oplus Q$ and the identifications
of its quotients with the child systems.  The next lemma recovers both.
Its second part concerns a further loss of information at a vertex:
a coarse quotient is identified with a space $W$ carrying an internal
flag or decomposition, whereas an adjoining interval retains only the
resulting finer flag and its fine quotient identifications.

\begin{lemma}\label{lem:filtered-lift}
Let $U,Q$ be vector spaces, with flags $U_i,Q_i$, and let
\[
 0\longrightarrow U_i/U_{i-1}\longrightarrow W_i
 \longrightarrow Q_i/Q_{i-1}\longrightarrow0
 \quad(1\leq i\leq m)
\]
be exact sequences.  In $V=U\oplus Q$, consider flags $V_i$ inducing
the given flags on $U,Q$, together with compatible identifications
$V_i/V_{i-1}\simeq W_i$.  Their space is a nonempty homogeneous space
for the additive group $H=\Hom(Q,U)$, whose stabilizer is
\[
 K=\{h\in H:h(Q_i)\subset U_{i-1}\text{ for every }i\}.
\]
If a $W_i$ is further equipped with a flag or a direct-sum
decomposition, transport that structure to the lifted quotient,
then retain its induced refined flags and graded identifications
while forgetting the coarse identification with $W_i$.
Let $K_{\rm coarse}$ be the stabilizer in $H$ of the coarse flag
with its full identified quotients, and let $K_{\rm interval}$ be
the stabilizer of the induced refined flag with its identified
fine quotients.  They are given by the displayed stabilizer formula
for the respective flags.  For compatible coarse and refined data,
restriction gives a surjection
$H/K_{\rm coarse}\to H/K_{\rm interval}$, with affine fiber
$K_{\rm interval}/K_{\rm coarse}$.

These assertions hold over $R$ when all named graded modules are free.
After free splittings, all choices and all fibers have free integral
parameters.  In fixed-rank algebraic families the maps are locally
affine bundles.
\end{lemma}

\begin{proof}
Split the flags on $U,Q$ and the exact sequences for $W_i$.  Their
direct sum supplies a lift.  For any other lift, successively lift
basis vectors of $W_i$ to $V_i$, using the chosen lifts in $U$ on
$U_i/U_{i-1}$.  Adjust the remaining lifts by $V_{i-1}$ to give them
the chosen projections to $Q$.  This splits $V$ as the sum of its
$W_i$ in a way compatible with the fixed splittings on $U,Q$.
Two resulting isomorphisms differ by
\[
 (u,q)\longmapsto(u+h(q),q),\qquad h\in H,
\]
which proves transitivity.  Such an automorphism fixes a lifted flag
and its identified quotients exactly when $h(Q_i)\subset U_{i-1}$.

An internal flag or decomposition of a $W_i$ is restored through its
identified quotient.  An automorphism fixing that full identification
fixes all the induced fine data.  Forgetting the full identification
therefore enlarges the stabilizer:
$K_{\rm coarse}\subset K_{\rm interval}$.  Restriction is the
surjection between the corresponding homogeneous spaces.  Its fibers are
$K_{\rm interval}/K_{\rm coarse}$.  In splittings compatible with
the given refinements these stabilizers are sums of matrix blocks,
and the indicated inclusion is split.  This proves the assertions
about integral parameters and, after local splittings of vector
bundles, the assertion for families.
\end{proof}

For example, let $U,Q$ be lines and take $W=U\oplus Q$ with its
standard exact sequence.  An identification of $V$ with $W$ inducing
the identity on $U,Q$ can vary by any element of $\Hom(Q,U)$.  If we
retain only the flag $0\subset U\subset V$ and its two identified
quotients, all those choices give the same data.  The coarse stabilizer
is zero and the refined-interval stabilizer is all of $\Hom(Q,U)$.
The reversal comes from forgetting the full coarse identification,
not from imposing fewer subspaces.

At a vertex there will be two adjoining intervals.  We use the
surjection to either one separately: an incoming lift extends to a
vertex lift, which then determines an outgoing lift.  We do not
prescribe two interval lifts independently or assert that every such
pair comes from a vertex lift.
